% ============================================================
% 片段：s05b  对应原文 section3.tex 第 308--408 行；PDF 第 58--61 页
% 内容：§5 Higher Group Symmetry（5.3 交叉扩张、5.4 Postnikov 类及两个例子）
% 译者：t-s05b   审校状态：初稿
% ============================================================

% ==== 原文 PDF 第 58 页（续）====
\subsection{交叉扩张}\label{sec:cross-extension}

本节我们试图在交叉模内部导出某种等价关系，并由此得到“弱”2-群\footnote{多数文献先引入弱 2-群，因为弱 2-群承载物理信息；但在我们的形式体系中，严格 2-群更自然地可以定义。}。

注意，交叉模$(G,H,\ptl , \rhd)$中群同态$\ptl$的存在会给出一个正合列
\be 1\to {\rm ker}\ptl \stack{i}{\longrightarrow} H \stack{\ptl}{\longrightarrow} G  \stack{p}{\longrightarrow} {\rm coker}\ptl \to 1 \ee
记$\Pi_2 \equiv {\rm ker}\ptl, \  \Pi_1\equiv {\rm coker}\ptl = G/ {\rm im}\ptl  $，这样做是合法的，因为$g(\ptl h)g^{-1} = \ptl(g\rhd h)\in {\rm im}\ptl$（${\rm im}\ptl$是正规子群）。这个正合列称为\textbf{交叉扩张}。

% ==== 原文 PDF 第 59 页 ====
现在可以如下定义一种等价关系：若图
\be
\label{2-group-equiv}
\begin{tikzcd}
	&& {H_1} & {G_1} \\
	1 & {\Pi_2} & {} && {\Pi_1} & 1 \\
	&& {H_2} & {G_2}
	\arrow[from=2-1, to=2-2]
	\arrow["{i_1}", from=2-2, to=1-3]
	\arrow["{i_2}"', from=2-2, to=3-3]
	\arrow["{\partial_1}",from=1-3, to=1-4]
	\arrow["{\partial_2}",from=3-3, to=3-4]
	\arrow["{t_H}"', from=1-3, to=3-3]
	\arrow["{t_G}"', from=1-4, to=3-4]
	\arrow["p_1",from=1-4, to=2-5]
	\arrow["p_2",from=3-4, to=2-5]
	\arrow[from=2-5, to=2-6]
\end{tikzcd}
\ee
交换，且与两个作用$\rhd_1, \ \rhd_2 $相容，即
\be t_H(g\rhd h) = t_G(g)\rhd' t_H(h) , \ \forall g\in G_1,\ H\in H_1  \ee
则称两个交叉模$(G_1,H_1,\ptl_1 , \rhd_2)$与$(G_2,H_2,\ptl_2 , \rhd_2)$弱等价。

这一等价类可以由$(\Pi_1,\Pi_2,\alpha,\beta)$很好地分类，其中$\Pi_{1,2}$即前面提到的两个群，$\alpha$是由$\rhd$诱导的作用$\alpha : \Pi_1\to {\rm Aut}(\Pi_2)$，而最值得注意的是，$\beta\in H^3_{\rm grp}(\Pi_1,\Pi_2)$是一个群上同调信息，称为\textbf{Postnikov 类}。

\begin{yizhu}这里的“弱等价”实为交叉扩张之间的同构：图 (5.7) 的两行分别是 $\Pi_2\to H_i\to G_i\to\Pi_1$ 两个交叉扩张，竖直箭头 $t_H,t_G$ 与 $i,\ptl,p$ 逐项相容，(5.8) 则是 $t_H$ 关于两个作用的等变性；二元组 $(t_H,t_G)$ 就是交叉模的同构。等价类由 $(\Pi_1,\Pi_2,\alpha,\beta)$ 分类这一点值得留意：前三个量只是“端点”数据（$\alpha$ 由 $\rhd$ 诱导），真正的扩张信息全在 $\beta\in H^3_{\rm grp}(\Pi_1,\Pi_2)$ 中——即使 $(\Pi_1,\Pi_2,\alpha)$ 相同，不同的交叉扩张仍可由 $\beta$ 区分。又，原文此处把第一个交叉模的作用指标写成 $\rhd_2$（应为 $\rhd_1$），(5.8) 的量词也写成 $H\in H_1$（应为 $h\in H_1$）；均照原文保留。\end{yizhu}

\subsection{Postnikov 类}\label{sec:Postnikov}

如何理解和计算 Postnikov 类这一信息？注意我们在第\ref{subsubsec:0+1D}节已提到群上同调。在那里，群上同调是衡量射影表示中结合律破坏程度的一个元素。正如我们即将看到的，Postnikov 类将扮演类似的角色：它分类交叉模中群作用的“扭曲”。

给定交叉扩张
\be 1\to \Pi_2 \stack{i}{\longrightarrow} H \stack{\ptl}{\longrightarrow} G  \stack{p}{\longrightarrow} \Pi_1 \to 1\,, \ee
我们按以下步骤计算 Postnikov 类：
\begin{enumerate}
    \item 选取提升映射$s: \Pi_1\to G$，使得$p\circ s = {\rm id}_{\Pi_1}$。注意它不必是同态，因为我们要把$f: \Pi_1\times \Pi_1 \to G$定义为
    \be  s(g)s(h) = s(g,h)f(g,h) \ . \ee
    此外，函数$f(g,h)$满足 3-上闭链条件
    \be
    s(g)f(h,k)s(g)^{-1}f(g,hk)=f(g,h)f(gh,k)\,.
    \ee
    \item 现在把$f$提升为函数$F:\Pi_1\times\Pi_1\rightarrow H$，要求对所有$g\in\Pi_1$有$\ptl (F(g))\equiv f(g)$。
    \item $F(g,h)$对（3-）上闭链条件的破坏可显式写为
    \be
    \label{cocycle-F}
    (s(g)\rhd F(h,k))F(g,hk)=i(c(g,h,k))F(g,h)F(gh,k)\,,
    \ee
    其中$c:\Pi_1\times\Pi_1\times\Pi_1\rightarrow\Pi_2\in H^3(B\Pi_1,\Pi_2) $就是我们想要的 Postnikov 类。
\end{enumerate}

\begin{yizhu}构造的要点是两次“提升”：先取集合层面的截面 $s$（它未必是同态），$f(g,h)$ 度量 $s$ 偏离同态的程度，而 (5.11) 表明 $f$ 满足 3-上闭链条件；再把 $f$ 提升为 $H$ 值函数 $F$。此时 $F$ 一般不再满足严格的上闭链条件，(5.12) 右端的 $i(c)$ 正是这一破坏，故 $c$ 的类 $[c]\in H^3(B\Pi_1,\Pi_2)$ 是“无法整体提升为严格上闭链”的障碍。它与提升的选取无关：若换用另一个截面 $s'$（它与 $s$ 相差一个 ${\rm im}\ptl$ 值函数），则 $f$ 与 $c$ 都只改变上边缘，因而 $[c]$ 只依赖交叉扩张本身。另，(5.10) 左端的 $s(g,h)$ 依 (5.11) 的用法应为 $s(gh)$，此处照原文保留。\end{yizhu}

% ==== 原文 PDF 第 60 页 ====
\subsubsection{例：$G = H = \mathbb{Z}_4$}
这是最简单的非平凡例子之一，其交叉扩张显式写为
\be 1\to \bbZ_2 \stack{\times 2}{\longrightarrow} \bbZ_4 \stack{\times 2}{\longrightarrow} \bbZ_4  \stack{\text{mod}\ 2}{\longrightarrow} \bbZ_2 \to 1 \ee
注意我们用$\{0,1,\dots,N-1\}$表示加法群$\mb{Z}_N$的元素。

由于${\rm Aut}(\bbZ_4) \cong \bbZ_2$，其中非平凡元素$\alpha$代表作用$(1,3)$（交换 1 与 3，其余元素不动）。此处的群上同调为
\be H^3_{\rm grp} (B\bbZ_2,\bbZ_2)=\bbZ_2 \ee
若取作用$\rhd = {\rm id}\in {\rm Aut}(\bbZ_4)$，则 Postnikov 类为$0\in H^3_{\rm grp} (B\bbZ_2,\bbZ_2)=\bbZ_2$；若取非平凡作用$\rhd = \alpha $，则 Postnikov 类为非平凡元素$1\in H^3_{\rm grp} (B\bbZ_2,\bbZ_2)$。

\begin{yizhu}此例的关键在于：$\Pi_2={\rm ker}(\times2)=\{0,2\}\cong\bbZ_2$，$\Pi_1=\bbZ_4/{\rm im}(\times2)\cong\bbZ_2$；非平凡自同构 $\alpha$ 即 $h\mapsto -h$（交换 $1$ 与 $3$），它限制在 $\Pi_2$ 上是恒等，因此无论 $\rhd$ 取 ${\rm id}$ 还是 $\alpha$，诱导的 $\alpha:\Pi_1\to{\rm Aut}(\Pi_2)$ 都平凡——两种取法在 $(\Pi_1,\Pi_2,\alpha)$ 层面完全不可区分，唯一的差别只能体现在 $\beta\in H^3_{\rm grp}(\bbZ_2,\bbZ_2)=\bbZ_2$ 上，这正是原文给出的 $0$ 与 $1$ 两个取值的来源；换言之，此例是 Postnikov 类作为独立数据的最小演示。具体计算按 (5.10)--(5.12) 进行：取提升 $s(\bar{1})=1$，由 (5.10) 得 $f$，再按 (5.11) 核对上闭链性质并把它提升为 $F$，最后由 (5.12) 读出 $c$；由于 $H^3_{\rm grp}(\bbZ_2,\bbZ_2)$ 只有两个元素，只需判断 $c$ 是否为零。\end{yizhu}

\subsubsection{例：$G = SU(2), \ H = \mathbb{Z}_{2l}$}
把注意力转向李群时，我们将看到丰富得多的结构。此时的交叉扩张为
\be 1\to \bbZ_{l} \stack{\times 2}{\longrightarrow} \bbZ_{2l} \stack{ {\rm mod} 2 \cdot I }{\longrightarrow} SU(2)  \stack{p}{\longrightarrow} SO(3) \to 1 \ee
我们希望把这个又长又复杂的正合列拆成若干我们能理解的短正合列，而这是可以做到的！
\be
    \begin{tikzcd}
	{1}\arrow{r} & {\mb{Z}_{l}}\arrow{r}{i}\arrow{d} & {\mb{Z}_{2l}}\arrow{r}\arrow{d}{\ptl = {\rm mod} 2 \cdot I} & {\mb{Z}_2 = {\rm im}\ptl}\arrow{r}\arrow{d} & {1} \\
	{1}\arrow{r} & {\mb{Z}_2 ={\rm im}\ptl}\arrow{r} & {SU(2)}\arrow{r}{p} & {SO(3)}\arrow{r} & {1}
\end{tikzcd}
\ee
对任何交叉扩张，这一步骤都是良定义的。重要的信息如下：
\begin{itemize}
    \item 第一行
    \be\begin{tikzcd}
    \label{eq:1strow}
        {1}\arrow{r} & {\mb{Z}_{l}}\arrow{r}{i} & {\mb{Z}_{2l}}\arrow{r} & {\mb{Z}_2 = {\rm im}\ptl} \arrow{r} & {1}
    \end{tikzcd} \ee
    是${\rm im}\ptl$被$\Pi_2$的中心扩张，由$\omega_1 \in H^2_{\rm grp}( \bbZ_2, \bbZ_{2l} )$分类。
    \item 类似地，第二行
    \be\begin{tikzcd}
        {1}\arrow{r} & {\mb{Z}_2 ={\rm im}\ptl}\arrow{r} & {SU(2)}\arrow{r}{p} & {SO(3)}\arrow{r} & {1}
    \end{tikzcd}\ee
    也是中心扩张，由$\omega_2 \in H^2_{\rm grp}( \Pi_1, {\rm im}\ptl ) = H^2_{\rm grp}(SO(3),\bbZ_2)$分类，它描述把$SO(3)$-丛提升为$SU(2)$-丛的障碍，而这恰好是第二 Stiefel–Whitney 类。对其它连续的$\Pi_1$与$G$，丛提升的解释同样成立。
    % ==== 原文 PDF 第 61 页 ====
    \item 那么前面所说的 Postnikov 类呢？三阶 Postnikov 类与二阶 Stiefel–Whitney 类之间是什么关系？这与第一行有关。首先介绍同调论中的一个著名工具。
    \begin{theorem}[Bockstein]
        给定阿贝尔群的正合列
        \be 0\to P\to Q\to R\to 0  \ee
        存在上同调群的同态$B_i : H^i(C,R)\to H^{i+1}(C,P)$，从而有长正合列
        \be  \cdots \to H^{n}(C,P) \to H^{n}(C,Q) \to H^{n}(C,R) \to
 H^{n+1}(C,P)  \to H^{n+1}(C,Q)\to \cdots  \ee
    \end{theorem}
    由这一定理便清楚可见：由于式~(\ref{eq:1strow}) 是正合列，我们可以诱导出 Bockstein 同态
    \be  B: H^2_{\rm grp}(\Pi_1,{\rm im}\ptl) \to H^3_{\rm grp}(\Pi_1,\Pi_2) , \ \ \omega_2\mapsto \beta \ee
    它恰好就是 Postnikov 类。
\end{itemize}

\begin{yizhu}这一例把三类信息串了起来：第一行 (5.17) 是系数正合列 $1\to\Pi_2\to\bbZ_{2l}\to{\rm im}\ptl\to1$；第二行 (5.18) 的分类量 $\omega_2$ 是 2-上闭链，恰好是 $SO(3)$-丛提升为 $SU(2)$-丛的障碍，即第二 Stiefel–Whitney 类 $w_2$；而 Bockstein 同态 (5.21) 正是用第一行的系数正合列把 $H^2$ 中的 $\omega_2$ 送到 $H^3$ 中的 $\beta$，即 $B(\omega_2)=\beta$，因此三阶 Postnikov 类与二阶 Stiefel–Whitney 类是同一障碍在不同次数上的表现。$\omega_1\in H^2_{\rm grp}(\bbZ_2,\bbZ_{2l})$ 则分类第一行这个中心扩张（它给出 Bockstein 所需的系数正合列）。\end{yizhu}
